$M_{R}^{-1} = \alpha^{-1}\Pi + M_{3}^{-1}\J/3$ by Lemma \claimref{neutrino.invariant-symmetric-matrices}. Expand $\mD M_{R}^{-1}\mD^{\mathsf{T}} = d^{2}(\J + \delta A)(\alpha^{-1}\Pi + M_{3}^{-1}\J/3)(\J + \delta A^{\mathsf{T}})$. Since $\J\Pi = \Pi\J = 0$, the $\Pi$ term leaves only $\alpha^{-1}\delta^{2}A\Pi A^{\mathsf{T}}$; since $\J = u'u'^{\mathsf{T}}$, the $\J/3$ term is $M_{3}^{-1}\tfrac13\,(\J + \delta A)u'u'^{\mathsf{T}}(\J + \delta A^{\mathsf{T}}) = M_{3}^{-1}\tfrac13\,qq^{\mathsf{T}}$. The identity $\J X\J = (\sum X_{ij})\J$ is $u'^{\mathsf{T}}Xu' = \sum_{ij}X_{ij}$. The second term is rank one along $q = 3u' + O(\delta)$ with eigenvalue $d^{2}|q|^{2}/(3M_{3}) = 9d^{2}/M_{3}\,(1 + O(\delta))$; at fixed $\alpha$ and $M_{3}$ the first term is $O(\delta^{2})$ and the second $O(1)$, so to leading order in $\delta$ the second dominates, the top eigenvector of the sum is $q + O(\delta^{2}M_{3}/\alpha)$, and the remaining two eigenvalues are those of the first term compressed to $q^{\perp} = u'^{\perp} + O(\delta)$: that is, $(d^{2}\delta^{2}/\alpha)$ times the eigenvalues of $\Pi A\Pi A^{\mathsf{T}}\Pi$. (The range of $A\Pi A^{\mathsf{T}}$ is $A(u'^{\perp})$, which is not orthogonal to $u'$ unless $A^{\mathsf{T}}u' \parallel u'$, so the eigenvalues of the first term itself are not the light masses; for $A = e_{3}e_{3}^{\mathsf{T}}$ the two differ by the factor $\tfrac23$ exactly.) The script verifies the exact identity on rational instances of $(\alpha, \beta, d, \delta, A)$; for $A = e_{3}e_{3}^{\mathsf{T}}$ (the $S_{3} \to S_{2}$ direction, $\lambda_{A} = \tfrac49$, $\lambda'_{A} = 0$) it diagonalises the light matrix exactly, the characteristic polynomial factoring as $\lambda$ times a quadratic, and certifies $m_{2}/m_{3} \in (0.168, 0.172)$ at $\delta = 0.286$ and $M_{3}/\alpha = 66$ (the instance of Remark \ref{rem:branch}), against $0.267$ from the leading-order formula there; and, at the small-breaking end, that the exact ratio at $\delta = 0.006$ is within four percent of the $\lambda_{A} = \tfrac49$ formula and about two thirds of the $A\Pi A^{\mathsf{T}}$ one.
